\chapter{Atlas Data Services, Serverless Integration, and Reverse ETL}
\index{change streams}\index{CDC}\index{resume token}\index{Atlas Triggers}\index{Atlas Functions}\index{reverse ETL}\index{Data API}\index{GraphQL}\index{serverless}%

\begin{objectives}
\begin{itemize}
    \item Build resumable change-stream consumers with persisted tokens, pre/post-images, and operation filters.
    \item Trigger Atlas Functions on database writes and compose event-driven serverless flows.
    \item Implement reverse ETL: curated aggregates pushed back into operational systems (CRM sync) with identity resolution and sync-frequency trade-offs.
    \item Expose HTTPS-based database interactions through the Atlas Data API and GraphQL API.
    \item Size serverless connection pools against cold-start storms.
    \item Architect an event-driven serverless order-fulfillment platform in the Friday engagement.
\end{itemize}
\end{objectives}

\begin{crosscloud}
Change data capture, serverless triggers, and reverse ETL are the connective tissue of every modern data platform; each cloud has a native variant, and the discipline---checkpoint, idempotency, replay---is identical across all of them.

\vspace{2mm}
{\footnotesize
\begin{tabular}{@{}p{2.9cm}p{3.0cm}p{3.0cm}p{3.0cm}@{}} 
\toprule
\textbf{Concept} & \textbf{AWS} & \textbf{Azure} & \textbf{GCP} \\
\midrule
Database CDC & DynamoDB Streams / DMS CDC & Cosmos DB change feed & Firestore triggers / Datastream \\
Serverless trigger & Lambda on streams & Functions on change feed & Cloud Functions / Run \\
HTTPS data API & DynamoDB API / RDS Data API & Cosmos DB REST & Firestore REST \\
Reverse ETL & AppFlow / custom & Synapse Link + Logic Apps & Datastream + Workflows \\
Event backbone & Kinesis / EventBridge & Event Hubs / Event Grid & Pub/Sub \\
\addlinespace
Snowflake / Fabric & \multicolumn{3}{l}{Streams + Tasks for CDC; reverse ETL into operational stores via connectors; Fabric Activator for triggers} \\
\bottomrule
\end{tabular}}

\vspace{1mm}
MongoDB change streams are distinguished by what they ride on: the same oplog that powers replication, so CDC is a first-class, total-order-per-shard, transactional-integrity-preserving API rather than a bolt-on log scrape \cite{mongodbChangeStreamsDocs,mongodbOplogDocs}.
\end{crosscloud}

\section{Week 16 Roadmap: Data Services Around the Core}
Every write your applications make is already an event; this week turns that capability into serverless architecture. Monday masters the change-stream consumer (resume tokens, pre/post-images, filters); Tuesday adds Atlas Triggers and Functions and the reverse-ETL pattern; Wednesday covers HTTPS-based access via the Data API and GraphQL and the serverless connection-pooling discipline; Thursday builds the resumable microservice; Friday composes the event-driven fulfillment platform \cite{mongodbChangeStreamsDocs}.

\begin{definitionbox}
\textbf{Reverse ETL} is the movement of curated, warehouse- or lake-computed aggregates (customer-360 scores, propensity segments) back into operational systems---CRMs, ad platforms, support tooling---so that analytics changes front-line behavior. \textbf{Identity resolution} is the matching layer that binds operational records to analytical entities on stable keys or hashed identifiers.
\end{definitionbox}

\begin{keyterms}
\textbf{Resume token}: an opaque position marker; the consumer's checkpoint. \textbf{Pre/post-image}: the document before/after the change, when collection-level capture is enabled. \textbf{Atlas Trigger}: a database or scheduled event that invokes an Atlas Function. \textbf{Data API}: HTTPS endpoint for CRUD and aggregation without a driver. \textbf{Oplog window}: how far back tokens remain valid. \textbf{At-least-once}: the delivery guarantee that makes sink idempotency mandatory.
\end{keyterms}

\begin{definitionbox}
A \textbf{change stream} is a cursor over a real-time, ordered feed of data changes (inserts, updates, deletes, DDL) for a collection, database, or deployment, backed by the oplog and resumable from a persisted token.
\end{definitionbox}

\section{Monday: Change Stream Fundamentals}
\subsection{What Streams Over What}
Streams open at three scopes: one collection, one database, or the whole deployment. Events carry \texttt{operationType}, the document key, and---by configuration---the \texttt{fullDocument} (post-image lookup), pre-images, and update descriptions. Ordering is total per shard and per document; cross-shard ordering is causal, not global, which matters when a consumer assumes timeline order across tenants.

The filtering pipeline is pushed into the server: \texttt{db.orders.watch([\{ \$match: \{ "fullDocument.status": "paid" \} \}])} streams only matching events. Filter early for the same reason Chapter 10 filters early: events you do not consume still cost oplog reads and network.

\subsection{Resume Tokens and the Oplog Window}
Every event carries an \texttt{\_id} that is its resume token. Persisting tokens after successful processing makes consumers crash-safe: reopen with \texttt{resumeAfter} and continue exactly where you left off. The boundary condition is the oplog window (Chapter 5): a token whose entries have rolled off is dead. The re-bootstrap procedure is the syllabus's production note made concrete: stop the consumer, snapshot the source, open a fresh stream from the post-snapshot point, and resume token checkpointing from there. Size the oplog so your worst-case consumer downtime fits the window with margin \cite{mongodbOplogDocs}.

\begin{pitfall}
\textbf{Checkpointing before processing.} Saving the token before the sink write loses events on crash; saving it after risks duplicate processing on crash. The correct pair is at-least-once delivery plus an idempotent sink---there is no token-timing trick that yields exactly-once without idempotency.
\end{pitfall}

\subsection{Pre/Post-Images and Transactional Integrity}
Enabling \texttt{changeStreamPreAndPostImages} on a collection adds the before-image to update/delete events---the difference between ``something changed'' and ``field X went from A to B,'' which audit and downstream-diff consumers require. Multi-document transactions stream their operations together and in order at commit, preserving the Chapter 7 invariant for consumers: a capture-and-ledger-post transaction never appears half-applied in the stream.

\subsection{The Reference Consumer}
\begin{lstlisting}[language=Python, caption={Resumable change-stream consumer with post-processing checkpoint.}]
from pymongo import MongoClient

client = MongoClient("mongodb://localhost:27017/?replicaSet=labrs")
db = client.shop
pipeline = [{"$match": {"operationType": {"$in": ["insert", "update"]}}}]

def consume():
    token = load_resume_token(db)              # persisted checkpoint
    kwargs = {"pipeline": pipeline, "full_document": "updateLookup"}
    if token:
        kwargs["resume_after"] = token
    try:
        with db.orders.watch(**kwargs) as stream:
            for change in stream:
                process_idempotent(change)     # sink upsert by document key
                save_resume_token(db, stream.resume_token)  # AFTER processing
    except Exception:
        consume()                              # re-open from last token

def process_idempotent(change):
    key = change["documentKey"]["_id"]
    doc = change.get("fullDocument")
    db.search_index.replace_one({"_id": key}, {"_id": key, "doc": doc},
                                upsert=True)   # replay-safe
\end{lstlisting}

\subsection{Idempotent Sinks and Dead Letters}
At-least-once delivery is the honest model: crashes between sink write and token save replay events. Sinks must therefore be idempotent---upserts keyed on the document key or a deterministic event key---or carry a dedup store. Events that fail processing repeatedly (poison events: schema drift, oversized payloads) must not stall the stream forever; after a bounded retry count, route them to a dead-letter collection with the error attached, and alert on DLQ depth.

\begin{tipbox}
Include the resume token or a monotonic event key in the sink document. It makes ``did we process event X?'' answerable and turns replay debugging from archaeology into a query.
\end{tipbox}

\section{Tuesday: Atlas Triggers, Functions, and Reverse ETL}
\subsection{Database Triggers and Functions}
An Atlas Trigger fires an Atlas Function on database change events (itself a managed change stream) or on a schedule. Functions are serverless JavaScript executing beside the cluster: no driver management, no connection code, billed per execution. The canonical use is the reaction path: an order write triggers a function that enriches the document, calls an external API, or appends to a processing queue. The design limits mirror all serverless triggers: functions are best for short, stateless reactions; heavy computation belongs in Chapter 15's stream processors; and every function needs the same idempotency discipline, because trigger delivery is at-least-once.

\subsection{Reverse ETL: Closing the Analytics Loop}
Reverse ETL pushes curated aggregates---customer-360 scores, churn propensities, segment memberships---from the analytical layer back into operational systems. The MongoDB-native pattern: an aggregation computes the score documents (Chapters 10--11), a change stream on the results collection fires an Atlas Function that calls the CRM's API. The two hard parts are \textbf{identity resolution}---matching operational contacts to analytical entities on stable keys (customer ID) or hashed emails, with an explicit unmatched-record quarantine---and the \textbf{sync-frequency trade-off}: real-time pushes are fresh but hammer the CRM's rate limits and amplify churn (a recomputed score that oscillates by 0.01 spams the CRM nightly); batch syncs with a materiality threshold (``push only if the score moved $>5\%$'') are calmer and cheaper. Choose per destination, and document the freshness each downstream system can rely on.

\begin{pitfall}
\textbf{Reverse ETL without idempotency or materiality.} A function that PUTs every score change into the CRM on every pipeline run will rate-limit itself and corrupt audit trails downstream. Upsert by external key, push only material changes, and log every call with the source watermark.
\end{pitfall}

\section{Wednesday: Atlas Data API, GraphQL, and Serverless Pooling}
\subsection{HTTPS-Based Database Access}
The Atlas Data API exposes CRUD and aggregation over HTTPS---no driver, no persistent connection---for environments where drivers are impractical: edge workers, sandboxed automation, third-party integrations. A hosted GraphQL API (App Services lineage) adds typed schemas and client-generated queries over the same collections. Both trade driver-level control (sessions, transactions, fine-grained pooling) for ubiquity; the rule of thumb is Data API for integration glue and low-rate paths, drivers for anything latency- or throughput-sensitive.

\subsection{Serverless Connection Pooling}
Serverless functions connecting directly to Atlas hit a scaling wall: each cold-starting function instance opens a fresh connection pool, and a burst of hundreds of concurrent functions can exhaust \texttt{maxPoolSize} $\times$ instances against the cluster's connection limit. The syllabus's production note is the fix list: reuse the client across warm invocations (global/module scope outside the handler), set a small \texttt{maxPoolSize} (1--10) for short-lived functions, and prefer the Data API or a connection-pooling proxy for spiky workloads. Monitor connections as a percentage of maximum---the Chapter 23 alert at 80\% exists precisely for this failure.

\begin{tipbox}
A cold-start storm looks like a database outage from the application side (connection timeouts) and like a connection spike from the database side. The fix is on the client side: warm client reuse and bounded pools.
\end{tipbox}

\section{Thursday Hands-On Project: Search-Index Sync Microservice}
Build the resumable consumer from Monday as a real service: order changes flow through a change stream (and, for the light-reaction variant, an Atlas Function) into a derived search-index collection, surviving crashes mid-stream.

\subsection{Lab Protocol}
Seed the Chapter 5 replica set with an orders collection. Start the consumer, then run a writer loop that updates random orders. Mid-run: kill the consumer process (\texttt{Stop-Process -Force}), keep writing for 30 seconds, restart the consumer, and verify the search index converges to exactly the current order state with no missing and no duplicated updates. Then inject a poison event (an order with a field shape the sink rejects) and verify DLQ routing plus stream continuation.

\subsection*{Deliverables}
\begin{itemize}
    \item The consumer service with persisted resume tokens and idempotent sink.
    \item The crash-recovery test log showing convergence after restart.
    \item The poison-event test showing DLQ routing and uninterrupted processing.
\end{itemize}

\subsection*{Acceptance Criteria}
\begin{itemize}
    \item After the mid-stream crash, the search index equals the source state (verified by a diff script).
    \item Token persistence demonstrably occurs after sink writes (code review plus fault test).
    \item The poison event lands in the DLQ with its error, and the stream processes subsequent events.
\end{itemize}

\subsection*{Stretch Goals}
\begin{itemize}
    \item Let the consumer sleep past the oplog window and execute the full re-bootstrap procedure.
    \item Replace the sink with a \texttt{\$merge}-based materialized view refresh (Chapter 11) triggered by the stream.
\end{itemize}

\section{Friday Use Case and Architecture: Event-Driven Order Fulfillment}
\subsection{Scenario and Requirements}
An e-commerce platform must decouple order placement from fulfillment, inventory, notification, and analytics. Requirements: no lost events, no duplicate customer notifications, fulfillment reacting within seconds, and the ability to add new consumers without touching the order service.

\subsection{Architecture}
Order placement writes the order and an outbox document in one transaction (Chapter 7's machinery)---the \textbf{outbox pattern}, which eliminates the dual-write race (write the order, crash, never publish) that naive producers suffer: the business write and the event commit atomically, and the stream guarantees publication \cite{kleppmann2017ddia}. A change stream on \texttt{orders} (filtered to paid/status events) feeds Kafka via the source connector. Three independent consumer groups process events: fulfillment (exactly-once-effect via dedup-key upserts), notifications (idempotent by event key against a sent-messages store), and an ASP windowed analytics job (Chapter 15) maintaining five-minute sales aggregates merged into a dashboard collection. An Atlas Trigger on \texttt{orders} handles the light reactions (fraud-flag webhook, cache invalidation) that do not justify a consumer group. Resume tokens checkpoint per consumer; the DLQ fronts an on-call runbook; the oplog is sized for a 24-hour consumer outage. New consumers subscribe to the topic---the order service never learns they exist.

\begin{figure}[H]
    \centering
    \resizebox{0.98\textwidth}{!}{%
    \begin{tikzpicture}[
        box/.style={rectangle, rounded corners, draw=primary, fill=primary!6, minimum width=2.5cm, minimum height=1cm, align=center, font=\small\bfseries},
        bus/.style={rectangle, rounded corners, draw=codepurple, fill=codepurple!8, minimum width=2.5cm, minimum height=1cm, align=center, font=\small\bfseries},
        sink/.style={rectangle, rounded corners, draw=green!45!black, fill=green!8, minimum width=2.5cm, minimum height=1cm, align=center, font=\small\bfseries},
        arrow/.style={-{Stealth[length=2.5mm]}, draw=primary, thick}
    ]
        \node[box] (api) at (0,0) {Order API\\(txn: order +\\outbox)};
        \node[box] (cs) at (3.8,0) {Change\\Streams};
        \node[bus] (kafka) at (7.4,0) {Kafka\\topics};
        \node[sink] (ful) at (11.6,1.9) {Fulfillment\\(dedup upserts)};
        \node[sink] (notif) at (11.6,0) {Notifications\\(idempotent)};
        \node[sink] (asp) at (11.6,-1.9) {ASP windows\\$\to$ dashboard};
        \node[box] (dlq) at (7.4,-2.6) {DLQ + alerts};
        \draw[arrow] (api) -- (cs);
        \draw[arrow] (cs) -- (kafka);
        \draw[arrow] (kafka) -- (ful);
        \draw[arrow] (kafka) -- (notif);
        \draw[arrow] (kafka) -- (asp);
        \draw[arrow, dashed] (kafka) -- (dlq);
    \end{tikzpicture}%
    }
    \caption{Event-driven fulfillment: transactional outbox, change-stream capture, Kafka fan-out, idempotent sinks, dead-letter safety.}
    \label{fig:ch15-fulfillment}
\end{figure}

\subsection{Guardrails}
The operational contract: consumer lag is the paging metric (not error logs), token checkpoints are monitored for age, DLQ depth above zero pages during business hours, and a quarterly drill re-bootstraps one consumer group from snapshot to prove the runbook. Resume-token mortality is treated as a capacity fact: oplog window $\geq$ 2$\times$ the longest tolerable consumer downtime \cite{mongodbChangeStreamsDocs}.

\section{Interview Q\&A}
\subsection{CDC and Streaming}
\begin{enumerate}
    \item \textbf{Q: What does a change-stream resume token identify?}\\
    A: A position in the oplog. Persisted after processing, it lets a consumer resume exactly after the last handled event---valid only while its oplog entries survive.

    \item \textbf{Q: Why persist the resume token only after processing an event?}\\
    A: Saving it before the sink write loses events on crash; saving it after means a crash replays the event---at-least-once, which the idempotent sink absorbs.

    \item \textbf{Q: What is reverse ETL, and what problem does identity resolution solve in it?}\\
    A: Reverse ETL pushes curated analytical aggregates back into operational systems like CRMs. Identity resolution binds analytical entities to operational records on stable keys or hashed emails---without it, scores land on the wrong contact or duplicate.

    \item \textbf{Q: Real-time versus batch CRM sync?}\\
    A: Real-time is fresh but hits rate limits and pushes noise (oscillating scores); batch with a materiality threshold is calmer and cheaper. Choose per destination and document the freshness contract.

    \item \textbf{Q: When does a resume token die?}\\
    A: When its oplog entries roll past the oplog window---a consumer paused longer than the window must re-bootstrap from a snapshot.

    \item \textbf{Q: Why the outbox pattern instead of ``write then publish''?}\\
    A: The two-write sequence can crash in between, losing the event. The outbox commits the event with the business write in one transaction; the stream guarantees publication.

    \item \textbf{Q: How do serverless functions exhaust a cluster's connections?}\\
    A: Every cold-starting instance opens a fresh pool; hundreds of concurrent functions multiply \texttt{maxPoolSize} past the connection limit. Warm-scope the client, bound the pool, or use the Data API.
\end{enumerate}

\section{Further Reading}
MongoDB's change-stream documentation defines tokens, images, and scope \cite{mongodbChangeStreamsDocs}, with oplog sizing governing token lifetime \cite{mongodbOplogDocs}. Chapter 15 covers Atlas Stream Processing over Kafka sources and windowed sinks in depth \cite{mongodbStreamProcessingDocs}, and Atlas Functions, the Data API, and App Services are documented in the Atlas references \cite{mongodbAtlasMonitoringDocs}. Kleppmann's stream chapters are the canonical treatment of CDC, logs, and exactly-once effects \cite{kleppmann2017ddia}, and the Apache Kafka project documents the backbone semantics \cite{apacheKafkaRepo}.

\section*{Key Takeaways}
\begin{itemize}
    \item Change streams expose the oplog as an ordered, resumable, server-filtered event API.
    \item Persist resume tokens after processing; pair at-least-once delivery with idempotent sinks.
    \item Token lifetime equals the oplog window; re-bootstrap is snapshot plus fresh stream.
    \item Atlas Triggers and Functions handle short, stateless reactions; heavy computation belongs to stream processors (Chapter 15).
    \item Reverse ETL closes the analytics loop: push by stable identity keys, only on material change, with every call logged.
    \item Serverless clients reuse warm connections and bound pools---or use the Data API for spiky, low-rate paths.
    \item Consumer lag, token age, and DLQ depth are the platform's paging metrics.
\end{itemize}

\section{Exercises}
\begin{enumerate}
    \item \textbf{(Conceptual)} Why can cross-shard change streams not guarantee global time ordering?
    \item \textbf{(Conceptual)} Compare change streams with polling on an \texttt{updatedAt} field: list three failure modes of polling that streams eliminate.
    \item \textbf{(Conceptual)} When do you need pre-images, and what do they cost?
    \item \textbf{(Conceptual)} A reverse-ETL sync oscillates a customer's score between 0.71 and 0.72 nightly. What two controls stop the CRM spam?
    \item \textbf{(Hands-on)} Add pre/post-image capture to the lab and build an audit trail of balance changes.
    \item \textbf{(Hands-on)} Implement the outbox pattern on the lab cluster and verify atomic order+event commit under a primary kill.
    \item \textbf{(Design)} Design the CDC architecture for syncing MongoDB into Snowflake for analytics, including re-bootstrap and schema drift.
    \item \textbf{(Design)} Design a reverse-ETL sync from a customer-360 collection to a CRM: identity resolution, materiality threshold, rate limits, and the quarantine for unmatched records.
    \item \textbf{(Interview)} ``We saw duplicate notifications after every deploy.'' Explain the mechanism and the two-part fix.
\end{enumerate}

\section{Milestone 4 Capstone: Agentic E-Commerce Recommendation and Search Engine}\label{sec:ch16-capstone}

\begin{capstonebox}
\textbf{Mission:} build the Milestone-4 system: a product catalog in Atlas with scripted BM25 and vector indexes, hybrid search with measured relevance, and an Atlas Stream Processing (or change-stream) path that turns five-minute user-activity windows into personalized recommendation documents.

\textbf{Estimated time:} 6--8 hours.

\textbf{What you will practice:}
\begin{itemize}
    \item End-to-end RAG-grade retrieval: chunking, embedding, indexing, fusion, evaluation.
    \item Streaming windowed personalization with idempotent \texttt{\$merge} sinks.
    \item Index and embedding lifecycle as code with CI guards.
\end{itemize}
\end{capstonebox}

\subsection*{Scenario}
The retailer from earlier capstones wants conversational product discovery plus real-time personalization: users search in natural language, and the recommendation rail updates within minutes of browsing behavior.

\subsection*{Requirements}
\begin{itemize}
    \item Seed a catalog (${\geq}$50k products); generate 768-d embeddings with \texttt{model\_version}; script both index creations (no UI clicks).
    \item Hybrid retrieval endpoint with RRF and category/tenant pre-filters; recall@10 evaluated on a labeled set.
    \item Streaming path: clickstream events $\to$ five-minute windows $\to$ \texttt{\$merge} upserts keyed on \texttt{user\_id + window\_start} into \texttt{recommendations}.
    \item At least three data-quality tests (relevance spot-checks, embedding-dimension validation, windowed-aggregation correctness) with a failing case.
    \item Monitoring evidence (search index status, stream health) and a cost estimate with two documented optimizations.
\end{itemize}

\subsection*{Implementation Steps}
\begin{enumerate}
    \item \textbf{Seed and embed.} Script the catalog load and embedding pass with dimension validation.
    \item \textbf{Index.} Apply both index definitions via API/Terraform; poll index status before serving.
    \item \textbf{Retrieve.} Implement and evaluate hybrid search; tune \texttt{numCandidates} on the recall curve.
    \item \textbf{Stream.} Build the windowed personalization with deterministic merge keys; replay-test for duplicates.
    \item \textbf{Operate.} Wire the data-quality tests, monitoring evidence, and cost note.
\end{enumerate}

\subsection*{Deliverables}
\begin{itemize}
    \item All seeding, indexing, retrieval, and streaming code; the evaluation table; the test suite.
    \item Monitoring/cost evidence and the architecture diagram.
    \item A demo script showing natural-language search and a recommendation updating within one window.
\end{itemize}

\subsection*{Acceptance Criteria}
\begin{itemize}
    \item Hybrid recall@10 beats both single-path baselines on the labeled set.
    \item Replaying the event stream produces zero duplicate recommendation rows.
    \item The dimension-validation test fails loudly on a simulated model drift.
    \item The system rebuilds from the repo on a clean environment.
\end{itemize}

\subsection*{Stretch Goals}
\begin{itemize}
    \item Add a cross-encoder reranker and quantify the recall@5 lift.
    \item Implement the blue/green embedding upgrade on the live catalog with zero retrieval downtime.
\end{itemize}